\section{Introduction}\label{sec:introduction}

Recent years have witnessed sensational advances of   reinforcement learning (RL) in many prominent  sequential  decision-making problems, such as playing the game of Go \cite{silver2016mastering, silver2017mastering}, playing   real-time strategy games  \cite{OpenAI_dota,alphastarblog},   robotic control  \cite{kober2013reinforcement,lillicrap2016continuous}, playing card games \cite{brown2017libratus,brown2019superhuman}, and autonomous driving \cite{shalev2016safe}, especially 
accompanied with  the development of  deep neural networks (DNNs) for function approximation \cite{mnih2015human}.  
Intriguingly, most of the successful applications involve the participation of more than one single agent/player\footnote{Hereafter, we will use \emph{agent} and \emph{player} interchangeably. }, which should be modeled systematically as multi-agent RL (MARL) problems. Specifically, MARL  addresses the sequential decision-making problem of  multiple autonomous  agents that operate in a common environment, each of which aims to optimize its own long-term return by interacting with the environment and other agents \cite{bu2008comprehensive}.  Besides the aforementioned popular ones, learning in multi-agent systems finds potential  applications in other subareas, including cyber-physical systems \cite{adler2002cooperative,wang2016towards}, finance  \cite{lee2002stock,lee2007multiagent}, sensor/communication networks \cite{cortes2004coverage,choi2009distributed}, and social science \cite{castelfranchi2001theory,leibo2017multi}. 

 
 
Largely, MARL algorithms can be placed   into   three groups, \emph{fully cooperative}, \emph{fully competitive}, and \emph{a mix of the two},  depending on the types of settings they address. In particular, in  the cooperative setting, agents   collaborate to optimize a common long-term return; while in the competitive setting, the return of agents usually sum up to zero. The mixed setting involves both cooperative and competitive agents, with general-sum returns. Modeling disparate MARL settings requires frameworks spanning from optimization theory, dynamic programming, game theory, and decentralized control, see \S\ref{subsec:MARL_framework} for more detailed discussions. In spite of these   existing multiple frameworks, several challenges in MARL  are in fact common across the different  settings, especially for the theoretical analysis.  Specifically, first, the learning goals in MARL are \emph{multi-dimensional}, as the objectives of all agents are not necessarily aligned, which brings up the challenge of dealing with equilibrium points, as well as  some additional   performance criteria beyond return-optimization, such as  the efficiency of communication/coordination, and  robustness against potential adversarial agents. Moreover, as all agents are improving their policies according to their own interests   concurrently, the environment faced by each agent becomes \emph{non-stationary}. This breaks or invalidates the basic framework of most  theoretical analyses in the single-agent setting. Furthermore, the joint action space that increases exponentially with the number of agents may cause scalability issues, known as the \emph{combinatorial nature} of MARL \cite{hernandez2018multiagent}. Additionally, the information structure, i.e., \emph{who knows what},  in MARL is more involved, as each agent has limited access to the observations of others, leading to possibly suboptimal decision rules locally. A detailed elaboration on the underlying  challenges can be found in \S\ref{sec:challenges}. 


There has in fact  been no shortage of efforts attempting to address the above   challenges. See \cite{bu2008comprehensive} for a comprehensive overview of earlier theories and algorithms on MARL. Recently, this  domain has gained resurgence of interest due to the advances of single-agent RL techniques. Indeed, a huge volume of work on MARL has appeared   lately, focusing on either identifying new learning criteria and/or setups \cite{foerster2016learning,zazo2016dynamic,zhang2018fully,subramanian2019reinforcement}, or developing new algorithms for existing setups, thanks to the development of deep learning \cite{heinrich2016deep,lowe2017multi,foerster2017counterfactual,gupta2017cooperative,omidshafiei2017deep,kawamura2017neural,zhang2019monte}, operations  research \cite{mazumdar2018convergence,jin2019minmax,zhang2019policyb,sidford2019solving}, and multi-agent systems \cite{oliehoek2016concise,arslan2017decentralized,yongacoglu2019learning,zhang2019online}.  
Nevertheless, not all the efforts are placed under rigorous theoretical footings, partly due to the limited understanding of even single-agent deep RL theories, and partly due to the inherent challenges in  multi-agent settings. As a consequence, it is imperative to review and organize the MARL algorithms with theoretical guarantees, in order to highlight the boundary of existing research endeavors, and stimulate potential future directions on this topic.   
%The resurgence is also driven by several new application scenarios, such as Poker AI and 
 
In this chapter, we provide a selective overview of  theories and algorithms in MARL, together with several significant while challenging applications. More specifically, we focus on two representative frameworks of MARL, namely, Markov/stochastic games and extensive-form games, in discrete-time settings as in standard single-agent RL. In conformity with the aforementioned three groups, we review and pay particular attention to   MARL algorithms with convergence and complexity analysis, most of which are fairly recent. 
With this focus in mind, we note that our overview is by no means comprehensive. 
In fact, besides the classical reference  \cite{bu2008comprehensive}, 
there are several other reviews on MARL that have appeared recently, due to the resurgence of MARL \cite{hernandez2017survey,hernandez2018multiagent,nguyen2018deep,oroojlooyjadid2019review}. 
We would like to emphasize that these reviews provide views and taxonomies  that are complementary   to ours:  \cite{hernandez2017survey} surveys the works that are specifically devised to address \emph{opponent-induced non-stationarity}, one of the challenges we discuss in \S\ref{sec:challenges}; \cite{hernandez2018multiagent,nguyen2018deep} are relatively more comprehensive, but with the focal point on \emph{deep} MARL, a subarea with scarce theories thus far; \cite{oroojlooyjadid2019review}, on the other hand, focuses on  algorithms in the  \emph{cooperative} setting only, though the review within this setting is extensive. 

Finally, we spotlight several new angles and  taxonomies that are comparatively underexplored in the existing MARL reviews, primarily owing to our own research endeavors and interests. First, we discuss the framework of extensive-form  games in MARL, in addition to the conventional one of Markov games, or even simplified repeated games \cite{bu2008comprehensive,hernandez2017survey,hernandez2018multiagent}; second, we summarize the progresses of a recently boosting subarea: decentralized MARL with \emph{networked} agents, as an extrapolation of our early works on this \cite{zhang2018fully,zhang18cdc,zhang2018finite}; third, we bring about the \emph{mean-field} regime into MARL, as a remedy for the case with an extremely large  population of agents; fourth, we highlight some recent advances in optimization theory, which shed lights on the (non-)convergence of policy-based methods for MARL, especially zero-sum games. We have also reviewed some of the literature on MARL in partially observed settings, but without using deep RL as heuristic solutions.  
We expect these new angles to help identify fruitful future research directions, and more importantly, inspire researchers with interests in establishing   rigorous theoretical foundations  on MARL.   

%several new things motivated by our own research.  


%Thee analysis however has faced great technical challenges when XXX. Few  theoretical foundations XXXX.



%talk about the challenges in MARL theory (briefly mention our challenges section). 

%Mention the new things that we would like to emphasize, or relatively missing in the previous survey: networked MARL, partially observed setting, extensive-form games, policy-based methods for games, mean-field regime, etc.

%\kznote{Considering the very recent reviews on  ``deep'' multi-agent RL \cite{hernandez2018multiagent,nguyen2018deep}, and specifically on ``cooperative deep'' MARL \cite{oroojlooyjadid2019review}, we should be focusing on a slightly different perspective, i.e., we may want to focus more on the ``theory'' part. We will just mention deep MARL a little bit, and talk more in the applications. This would thus be orthogonal to the literature.}


%\noindent \issue{We will emphasize that our review is by no means ``comprehensive''.}

%\issue{Also note that our focus is on ``discrete-time'', RL-type algorithms.}


%\issue{Add a note, saying that we use \emph{agent} and \emph{player} interchangeably. }

\vspace{8pt}
{\noindent \bf Roadmap.}
\vspace{2pt}
The remainder of this chapter is organized as follows. In \S\ref{sec:background}, we introduce the background of MARL: standard algorithms for single-agent RL, and  the  frameworks of MARL. In \S\ref{sec:challenges}, we summarize several challenges in developing MARL theory, in addition to the single-agent counterparts. A series of MARL algorithms, mostly with theoretical guarantees, are reviewed and organized in \S\ref{sec:algorithms},  according to the types of tasks they address. In \S\ref{sec:applications}, we briefly introduce a few recent successes of MARL driven by the algorithms mentioned, followed by conclusions and several open  research directions outlined in \S\ref{sec:conclusion}. 



